% =====================================================================
% MODULE: L3.3 — CAMPOS VECTORIALES CONSERVATIVOS E INTEGRAL DE LÍNEA
%          DE UN CAMPO ESCALAR
% =====================================================================
% Subtemas: T3.5, T3.6 (según Constitución de lecciones Cálculo III)
% Fuente LPS: LPS_L3_3.pdf
% =====================================================================

\chapter{Campos Vectoriales Conservativos e Integral de Línea de un Campo Escalar}

En este módulo se exploran dos conceptos fundamentales del cálculo
vectorial. Primero, los \textbf{campos vectoriales} — funciones que
asignan un vector a cada punto del espacio— y en particular los
\textbf{campos conservativos}, que son aquellos que provienen del
gradiente de una función escalar. Estos campos gozan de una propiedad
notable: el trabajo realizado a lo largo de una trayectoria no depende
del camino recorrido, sino únicamente de los puntos inicial y final.

En segundo lugar, se introduce la \textbf{integral de línea de un campo
escalar}, una herramienta que generaliza la integral definida a
trayectorias curvas en el plano o en el espacio. Esta integral permite
acumular cantidades como masa, carga o temperatura a lo largo de un
camino, y su cálculo se reduce a una integral ordinaria mediante una
parametrización adecuada de la curva.

Ambos conceptos son pilares del cálculo multivariable y tienen
aplicaciones directas en física, ingeniería y geometría. La conexión
entre ellos aparece cuando el campo escalar es una función potencial:
entonces la integral de línea del campo vectorial gradiente se evalúa
simplemente como la diferencia de potencial entre los extremos.

\section{Campos Vectoriales: Gradiente y Conservativos}

\subsection*{Definición de Campo Vectorial}

Un \textbf{campo vectorial} en \(\mathbb{R}^3\) es una función
\(\mathbf{F}: D \subseteq \mathbb{R}^3 \to \mathbb{R}^3\) que asigna a
cada punto \((x,y,z)\) de su dominio un vector tridimensional. Se
escribe como
\[
\mathbf{F}(x,y,z) = M(x,y,z)\,\mathbf{i} + N(x,y,z)\,\mathbf{j}
+ P(x,y,z)\,\mathbf{k},
\]
donde \(M,N,P\) son funciones escalares llamadas \textbf{componentes}
del campo. En \(\mathbb{R}^2\), un campo vectorial tiene la forma
\[
\mathbf{F}(x,y) = M(x,y)\,\mathbf{i} + N(x,y)\,\mathbf{j}.
\]

\begin{rem}
Los campos vectoriales modelan magnitudes que tienen dirección y
sentido en cada punto del espacio. Ejemplos físicos incluyen:
\begin{itemize}
  \item el campo de velocidades de un fluido (cada partícula se mueve
  con la velocidad indicada por el campo),
  \item el campo gravitacional (la fuerza que experimenta una masa
  unitaria),
  \item el campo eléctrico (la fuerza por unidad de carga),
  \item el campo magnético.
\end{itemize}
\end{rem}

\begin{definition}[Continuidad y diferenciabilidad de un campo vectorial]
Un campo vectorial \(\mathbf{F} = M\mathbf{i}+N\mathbf{j}+P\mathbf{k}\)
es \textbf{continuo} en una región \(D\) si cada una de sus funciones
componentes \(M,N,P\) es continua en \(D\). Es \textbf{diferenciable}
si cada componente tiene derivadas parciales continuas en \(D\).
\end{definition}

\subsection*{Campo Vectorial Gradiente}

Dada una función escalar \(f: D \subseteq \mathbb{R}^n \to \mathbb{R}\)
diferenciable, su \textbf{gradiente} es el campo vectorial definido por
\[
\nabla f(x,y) = f_x(x,y)\,\mathbf{i} + f_y(x,y)\,\mathbf{j}
\quad \text{en } \mathbb{R}^2,
\]
y
\[
\nabla f(x,y,z) = f_x(x,y,z)\,\mathbf{i}
+ f_y(x,y,z)\,\mathbf{j}
+ f_z(x,y,z)\,\mathbf{k}
\quad \text{en } \mathbb{R}^3.
\]

\begin{rem}[Interpretación geométrica del gradiente]
El vector \(\nabla f(P)\) apunta en la dirección de \textbf{máximo
crecimiento} de \(f\) en el punto \(P\), y su norma
\(\|\nabla f(P)\|\) mide la tasa de cambio de \(f\) en esa dirección.
Además, \(\nabla f(P)\) es perpendicular a las curvas (o superficies)
de nivel de \(f\) que pasan por \(P\).
\end{rem}

\subsection*{Campos Vectoriales Conservativos}

Un campo vectorial \(\mathbf{F}\) se dice \textbf{conservativo} si
existe una función escalar \(\Phi\), llamada \textbf{función
potencial}, tal que
\[
\mathbf{F} = \nabla \Phi.
\]
Es decir, un campo conservativo es el gradiente de alguna función
escalar. En este caso, el campo se denomina \textbf{campo gradiente}.

\begin{theorem}[Caracterización de campos conservativos en el plano]
Sean \(M(x,y)\) y \(N(x,y)\) funciones con derivadas parciales continuas
en una región abierta \(R\). El campo vectorial
\(\mathbf{F}(x,y) = M(x,y)\mathbf{i} + N(x,y)\mathbf{j}\) es
conservativo si y solo si
\[
\frac{\partial N}{\partial x} = \frac{\partial M}{\partial y}
\quad \text{en } R.
\]
\end{theorem}

\begin{rem}
La condición \(\partial N/\partial x = \partial M/\partial y\) es
necesaria y suficiente para la existencia de función potencial en
regiones \emph{simplemente conexas} (es decir, sin agujeros). En
regiones con agujeros, la igualdad de derivadas cruzadas no basta:
puede haber campos con rotacional nulo que no sean conservativos.
\end{rem}

\begin{rem}[Propiedades de los campos conservativos]
Si \(\mathbf{F} = \nabla \Phi\) es conservativo, entonces:
\begin{enumerate}
  \item La integral de línea de \(\mathbf{F}\) entre dos puntos
  \emph{no depende del camino} seguido, sino solo de los puntos
  inicial y final.
  \item En regiones simplemente conexas, su rotacional es nulo:
  \[
  \nabla \times \mathbf{F} = \mathbf{0}.
  \]
  \item El trabajo realizado por \(\mathbf{F}\) a lo largo de una
  curva cerrada es cero.
\end{enumerate}
\end{rem}

\begin{example}[Verificación de campo conservativo en el plano]
Determine si el campo
\[
\mathbf{F}(x,y) = (2x y)\,\mathbf{i} + (x^2 + 3y^2)\,\mathbf{j}
\]
es conservativo. En caso afirmativo, halle su función potencial.
\end{example}

\begin{myproof}
\textbf{Paso 1. Identificar componentes.}
Aquí \(M(x,y) = 2xy\) y \(N(x,y) = x^2 + 3y^2\).

\textbf{Paso 2. Verificar condición de igualdad de derivadas cruzadas.}
Calculamos:
\[
\frac{\partial N}{\partial x} = 2x, \qquad
\frac{\partial M}{\partial y} = 2x.
\]
Como \(\partial N/\partial x = \partial M/\partial y\), el campo
cumple la condición necesaria en todo \(\mathbb{R}^2\), que es una
región simplemente conexa. Por tanto, \(\mathbf{F}\) es conservativo.

\textbf{Paso 3. Hallar la función potencial.}
Buscamos \(\Phi(x,y)\) tal que \(\nabla \Phi = \mathbf{F}\):
\[
\Phi_x = 2xy, \qquad \Phi_y = x^2 + 3y^2.
\]
Integramos \(\Phi_x\) respecto a \(x\):
\[
\Phi(x,y) = \int 2xy \, dx = x^2 y + g(y).
\]
Derivamos respecto a \(y\):
\[
\Phi_y = x^2 + g'(y).
\]
Igualamos con la segunda componente:
\[
x^2 + g'(y) = x^2 + 3y^2 \implies g'(y) = 3y^2.
\]
Integrando: \(g(y) = y^3 + C\). Por tanto:
\[
\Phi(x,y) = x^2 y + y^3 + C.
\]
Verificamos: \(\nabla \Phi = (2xy)\,\mathbf{i} + (x^2+3y^2)\,\mathbf{j}
= \mathbf{F}\). \(\boxed{\text{Es conservativo con } \Phi = x^2 y + y^3 + C}\).
\end{myproof}

\section{Integral de Línea de un Campo Escalar}

La integral de línea de un campo escalar generaliza la integral
definida a curvas. En lugar de integrar sobre un intervalo en el eje
\(x\), se integra sobre una curva \(C\) en el plano o en el espacio,
acumulando los valores de una función escalar \(f\) a lo largo de la
curva.

\subsection*{Definición mediante suma de Riemann}

Sea \(C\) una curva suave en \(\mathbb{R}^2\) parametrizada por
\(\mathbf{r}(t) = (x(t), y(t))\), con \(a \le t \le b\). Dividimos
el intervalo \([a,b]\) en \(n\) subintervalos, lo que induce una
partición de la curva en pequeños segmentos de longitud \(\Delta s_i\).
En cada segmento elegimos un punto \((x_i^*, y_i^*)\) y formamos la
suma
\[
\sum_{i=1}^n f(x_i^*, y_i^*) \,\Delta s_i.
\]
La \textbf{integral de línea de \(f\) a lo largo de \(C\)} es el
límite de estas sumas cuando la partición se refina:
\[
\int_C f(x,y)\,ds = \lim_{n\to\infty}
\sum_{i=1}^n f(x_i^*, y_i^*) \,\Delta s_i.
\]
En \(\mathbb{R}^3\), la definición es análoga:
\[
\int_C f(x,y,z)\,ds = \lim_{n\to\infty}
\sum_{i=1}^n f(x_i^*, y_i^*, z_i^*) \,\Delta s_i.
\]

\begin{rem}
La notación \(ds\) representa el \textbf{elemento de longitud de
arco}. Para una curva parametrizada \(\mathbf{r}(t)\),
\[
ds = \|\mathbf{r}'(t)\|\,dt.
\]
Esta relación es la clave para evaluar la integral de línea mediante
una integral definida ordinaria.
\end{rem}

\begin{pasos}[Evaluación de una integral de línea de campo escalar]
Para calcular \(\int_C f\,ds\), se sigue este protocolo:
\begin{enumerate}
  \item \textbf{Parametrizar la curva.} Obtener una parametrización
  \(\mathbf{r}(t) = (x(t), y(t), z(t))\), \(a \le t \le b\), que
  recorra la curva \(C\).
  \item \textbf{Calcular la rapidez.} Hallar
  \(\|\mathbf{r}'(t)\| = \sqrt{(x')^2 + (y')^2 + (z')^2}\).
  \item \textbf{Sustituir en el integrando.} Reemplazar \(x,y,z\) por
  \(x(t),y(t),z(t)\) en \(f\).
  \item \textbf{Integrar.} Evaluar la integral definida
  \[
  \int_a^b f(x(t),y(t),z(t)) \,\|\mathbf{r}'(t)\|\,dt.
  \]
\end{enumerate}
\end{pasos}

\begin{theorem}[Teorema de evaluación para integral de línea de campo escalar]
Sea \(C\) una curva suave en \(\mathbb{R}^2\) parametrizada por
\(\mathbf{r}(t) = (x(t), y(t))\), \(a \le t \le b\), con \(x(t)\) y
\(y(t)\) derivables con derivadas continuas. Sea \(f\) continua en
una región que contiene a \(C\). Entonces
\[
\int_C f(x,y)\,ds =
\int_a^b f(x(t), y(t)) \sqrt{[x'(t)]^2 + [y'(t)]^2}\,dt.
\]
Análogamente, en \(\mathbb{R}^3\):
\[
\int_C f(x,y,z)\,ds =
\int_a^b f(x(t), y(t), z(t)) \sqrt{[x'(t)]^2 + [y'(t)]^2 + [z'(t)]^2}\,dt.
\]
\end{theorem}

\begin{rem}[Curvas suaves por partes]
Muchas curvas no son completamente suaves, pero pueden descomponerse en
un número finito de curvas suaves \(C_1, C_2, \dots, C_n\). En tal caso,
\[
\int_C f\,ds = \int_{C_1} f\,ds + \int_{C_2} f\,ds + \cdots + \int_{C_n} f\,ds.
\]
Esta propiedad es consecuencia directa de la aditividad de la integral
definida.
\end{rem}

\begin{rem}[Independencia de la orientación para campos escalares]
A diferencia de las integrales de línea de campos vectoriales, la
integral de línea de un campo escalar \emph{no} cambia de signo al
recorrer la curva en sentido contrario:
\[
\int_{-C} f\,ds = \int_C f\,ds.
\]
Esto se debe a que \(ds\) es siempre positivo, independientemente de
la orientación.
\end{rem}

\begin{example}[Integral de línea de un campo escalar en el plano]
Calcule \(\displaystyle \int_C (x + y)\,ds\), donde \(C\) es el
segmento de recta desde \((0,0)\) hasta \((1,2)\).
\end{example}

\begin{myproof}
\textbf{Paso 1. Parametrizar la curva.}
El segmento se parametriza como
\[
\mathbf{r}(t) = (1-t)(0,0) + t(1,2) = (t, 2t), \qquad 0 \le t \le 1.
\]

\textbf{Paso 2. Calcular la rapidez.}
\[
\mathbf{r}'(t) = (1, 2) \implies
\|\mathbf{r}'(t)\| = \sqrt{1^2 + 2^2} = \sqrt{5}.
\]

\textbf{Paso 3. Sustituir en el integrando.}
\(f(x,y) = x + y\), luego
\[
f(x(t), y(t)) = t + 2t = 3t.
\]

\textbf{Paso 4. Integrar.}
\[
\int_C (x+y)\,ds =
\int_0^1 3t \cdot \sqrt{5}\,dt =
3\sqrt{5} \int_0^1 t\,dt =
3\sqrt{5} \cdot \frac{1}{2}
= \boxed{\frac{3\sqrt{5}}{2}}.
\]
\end{myproof}

\begin{example}[Integral de línea en el espacio]
Calcule \(\displaystyle \int_C (x^2 + y^2 + z^2)\,ds\), donde \(C\) es
la hélice circular \(\mathbf{r}(t) = (\cos t, \sin t, t)\),
\(0 \le t \le 2\pi\).
\end{example}

\begin{myproof}
\textbf{Paso 1. Parametrización dada.}
\(x(t) = \cos t,\; y(t) = \sin t,\; z(t) = t\).

\textbf{Paso 2. Rapidez.}
\[
\mathbf{r}'(t) = (-\sin t, \cos t, 1) \implies
\|\mathbf{r}'(t)\| = \sqrt{\sin^2 t + \cos^2 t + 1} = \sqrt{2}.
\]

\textbf{Paso 3. Sustitución.}
\[
f(x,y,z) = x^2 + y^2 + z^2 = \cos^2 t + \sin^2 t + t^2 = 1 + t^2.
\]

\textbf{Paso 4. Integrar.}
\[
\int_C (x^2 + y^2 + z^2)\,ds =
\int_0^{2\pi} (1 + t^2)\sqrt{2}\,dt =
\sqrt{2} \left[ t + \frac{t^3}{3} \right]_0^{2\pi}
= \sqrt{2} \left( 2\pi + \frac{8\pi^3}{3} \right)
= \boxed{2\sqrt{2}\,\pi \left(1 + \frac{4\pi^2}{3}\right)}.
\]
\end{myproof}

\begin{example}[Integral sobre una curva suave por partes]
Calcule \(\displaystyle \int_C (x^2 + y^2)\,ds\), donde \(C\) es la
frontera del cuadrado con vértices \((0,0)\), \((1,0)\), \((1,1)\),
\((0,1)\), recorrida en sentido antihorario.
\end{example}

\begin{myproof}
La curva \(C\) es la unión de cuatro segmentos:
\(C_1: (0,0)\to(1,0)\), \(C_2: (1,0)\to(1,1)\),
\(C_3: (1,1)\to(0,1)\), \(C_4: (0,1)\to(0,0)\).

Por aditividad:
\[
\int_C f\,ds = \int_{C_1} f\,ds + \int_{C_2} f\,ds
+ \int_{C_3} f\,ds + \int_{C_4} f\,ds.
\]

\textbf{Segmento \(C_1\):} \((t,0)\), \(0\le t\le 1\).
\(\mathbf{r}'=(1,0)\), \(\|\mathbf{r}'\|=1\), \(f=t^2\).
\[
\int_{C_1} = \int_0^1 t^2\,dt = \frac{1}{3}.
\]

\textbf{Segmento \(C_2\):} \((1,t)\), \(0\le t\le 1\).
\(\|\mathbf{r}'\|=1\), \(f=1+t^2\).
\[
\int_{C_2} = \int_0^1 (1+t^2)\,dt = 1 + \frac{1}{3} = \frac{4}{3}.
\]

\textbf{Segmento \(C_3\):} \((1-t,1)\), \(0\le t\le 1\).
\(\|\mathbf{r}'\|=1\), \(f=(1-t)^2+1\).
\[
\int_{C_3} = \int_0^1 [(1-t)^2 + 1]\,dt
= \int_0^1 (t^2 - 2t + 2)\,dt
= \frac{1}{3} - 1 + 2 = \frac{4}{3}.
\]

\textbf{Segmento \(C_4\):} \((0,1-t)\), \(0\le t\le 1\).
\(\|\mathbf{r}'\|=1\), \(f=(1-t)^2\).
\[
\int_{C_4} = \int_0^1 (1-t)^2\,dt = \frac{1}{3}.
\]

Sumando:
\[
\int_C f\,ds = \frac{1}{3} + \frac{4}{3} + \frac{4}{3} + \frac{1}{3}
= \frac{10}{3}.
\]
\[
\boxed{\int_C (x^2+y^2)\,ds = \frac{10}{3}}.
\]
\end{myproof}

\section*{Conclusión}

En esta lección se han estudiado dos conceptos fundamentales:

\begin{enumerate}
  \item Los \textbf{campos vectoriales conservativos} como aquellos que
  provienen de un gradiente, caracterizados por la independencia de la
  trayectoria en la integral de línea y por la condición de igualdad
  de derivadas cruzadas (en el plano). Se introdujo la función
  potencial como el objeto central que permite evaluar el trabajo sin
  integrar a lo largo de la curva.

  \item La \textbf{integral de línea de un campo escalar}, que
  generaliza la integral definida a curvas en el plano y el espacio.
  Su cálculo se reduce a una integral ordinaria mediante una
  parametrización, y es independiente de la orientación de la curva.
\end{enumerate}

Estas herramientas son la antesala de conceptos más avanzados como la
integral de línea de campos vectoriales, el Teorema de Green y el
Teorema Fundamental para integrales de línea, que se abordarán en
lecciones posteriores.

\section*{Problemas Propuestos}

\begin{prob}[Campos vectoriales conservativos]
Determine si los siguientes campos son conservativos. En caso
afirmativo, halle una función potencial.
\begin{enumerate}[a)]
  \item \(\mathbf{F}(x,y) = (e^x \cos y)\,\mathbf{i} - (e^x \sin y)\,\mathbf{j}\).
  \item \(\mathbf{F}(x,y) = (y^2 + 2xy)\,\mathbf{i} + (x^2 + 2xy)\,\mathbf{j}\).
  \item \(\mathbf{F}(x,y,z) = (yz, xz, xy)\).
  \item \(\mathbf{F}(x,y) = \left(\dfrac{-y}{x^2+y^2},
  \dfrac{x}{x^2+y^2}\right)\) en la región \(\mathbb{R}^2 \setminus \{(0,0)\}\).
\end{enumerate}
\end{prob}

\begin{prob}[Cálculo de integrales de línea de campo escalar]
Evalúe las siguientes integrales de línea:
\begin{enumerate}[a)]
  \item \(\displaystyle \int_C (2x + 3y)\,ds\), donde \(C\) es el
  segmento de \((1,2)\) a \((4,6)\).
  \item \(\displaystyle \int_C (x^2 + y^2 + z^2)\,ds\), donde \(C\) es
  la curva \(\mathbf{r}(t) = (t, t^2, t^3)\), \(0\le t\le 1\).
  \item \(\displaystyle \int_C x\,ds\), donde \(C\) es el arco de
  la parábola \(y = x^2\) desde \((0,0)\) hasta \((2,4)\).
  \item \(\displaystyle \int_C \sqrt{1 + 4x^2}\,ds\), donde \(C\) es
  la curva \(y = x^2\), \(0\le x\le 1\). (Observe la relación con
  la longitud de arco).
\end{enumerate}
\end{prob}

\begin{prob}[Aplicación: masa de un alambre]
Un alambre tiene la forma de la hélice \(\mathbf{r}(t) = (2\cos t,
2\sin t, t)\), \(0 \le t \le 4\pi\). Su densidad lineal en cada punto
es \(\rho(x,y,z) = x^2 + y^2 + z^2\). Calcule la masa total del alambre.
\end{prob}

\begin{prob}[Curvas suaves por partes]
Calcule \(\displaystyle \int_C (xy)\,ds\), donde \(C\) es la frontera
del triángulo con vértices \((0,0)\), \((2,0)\), \((0,1)\), recorrida
en sentido antihorario.
\end{prob}

\begin{prob}[Análisis conceptual]
Responda con verdadero o falso. Justifique.
\begin{enumerate}[a)]
  \item Si \(\mathbf{F}(x,y) = \nabla f(x,y)\), entonces
  \(\mathbf{F}\) es conservativo.
  \item Si \(\partial N/\partial x = \partial M/\partial y\) en
  \(\mathbb{R}^2\), entonces \(\mathbf{F} = (M,N)\) es
  automáticamente conservativo.
  \item La integral de línea de un campo escalar cambia de signo
  si se recorre la curva en sentido contrario.
  \item Para una curva \(C\) suave por partes, la integral de línea
  se calcula como la suma de las integrales en cada tramo suave.
  \item El campo \(\mathbf{F}(x,y) = (y, -x)\) es conservativo en
  \(\mathbb{R}^2\).
\end{enumerate}
\end{prob}

\begin{prob}[Relación entre gradiente e integral de línea]
Sea \(f(x,y) = x^2 y + e^x \cos y\). Calcule \(\nabla f\) y evalúe
\(\displaystyle \int_C \nabla f \cdot d\mathbf{r}\) a lo largo de
cualquier curva \(C\) que vaya de \((0,0)\) a \((1,\pi/2)\).
¿Qué propiedad de los campos conservativos se está utilizando?
\end{prob}

\begin{prob}[Caso especial: integral de longitud de arco]
Demuestre que si \(f(x,y,z) = 1\), entonces
\[
\int_C 1\,ds = \text{longitud de arco de } C.
\]
Utilice esta observación para calcular la longitud de la curva
\(\mathbf{r}(t) = (t, t^2, t^3)\), \(0\le t\le 1\).
\end{prob}